% 04_gpu.tex -- GPU top level, one shader core, SIMT execution, latency hiding, comparison.
% Build: tectonic 04_gpu.tex
\documentclass[tikz,border=10pt,12pt]{standalone}
% soc_style.tex -- shared look for the SoC figures (same palette as the slide decks and the
% core register placement figure). Every figure does % soc_style.tex -- shared look for the SoC figures (same palette as the slide decks and the
% core register placement figure). Every figure does % soc_style.tex -- shared look for the SoC figures (same palette as the slide decks and the
% core register placement figure). Every figure does \input{soc_style}.
\usepackage[T1]{fontenc}
\usepackage[scaled=0.95]{helvet}
\renewcommand{\familydefault}{\sfdefault}
\usepackage{booktabs}
\usepackage{amsmath}
\usepackage{amssymb}
\usetikzlibrary{arrows.meta,calc,positioning,fit,shapes.geometric,shapes.misc}

% ---------------- palette ----------------
\definecolor{pcb}{HTML}{0B2E33}
\definecolor{pcbmid}{HTML}{1F5C61}
\definecolor{tint}{HTML}{E6F0EF}
\definecolor{copper}{HTML}{B5651D}
\definecolor{coppert}{HTML}{F6E8DC}
\definecolor{ink}{HTML}{1F2A2E}
\definecolor{muted}{HTML}{5B6B70}
\definecolor{hardf}{HTML}{DCE6F2}
\definecolor{harde}{HTML}{2F4A7A}
\definecolor{good}{HTML}{2E7D32}
\definecolor{bad}{HTML}{B3261E}
% one colour per kind of block: border (d) / fill (f)
\definecolor{cpud}{HTML}{2B5C9E} \definecolor{cpuf}{HTML}{D6E4F5}
\definecolor{gpud}{HTML}{2E7D32} \definecolor{gpuf}{HTML}{D5EDD0}
\definecolor{npud}{HTML}{0B2E33} \definecolor{npuf}{HTML}{BFE0DB}
\definecolor{memd}{HTML}{5E3A8C} \definecolor{memf}{HTML}{E4DAF3}
\definecolor{busd}{HTML}{B5651D} \definecolor{busf}{HTML}{F6E8DC}
\definecolor{prd}{HTML}{546E7A}  \definecolor{prf}{HTML}{E3EAED}
\definecolor{secd}{HTML}{B3261E} \definecolor{secf}{HTML}{F7D9D6}
\definecolor{pwrd}{HTML}{8A6800} \definecolor{pwrf}{HTML}{FFF0B8}
\definecolor{medd}{HTML}{8E2463} \definecolor{medf}{HTML}{F5D6E6}
\definecolor{skd}{HTML}{2E7D32}  \definecolor{skf}{HTML}{BFE3B4}   % skew buffers, as in the register figure

\tikzset{
  blk/.style 2 args={draw=#1, fill=#2, rounded corners=4pt, thick, align=center,
                     font=\scriptsize, inner sep=3pt},
  cpu/.style={blk={cpud}{cpuf}},  gpu/.style={blk={gpud}{gpuf}},
  npu/.style={blk={npud}{npuf}},  mem/.style={blk={memd}{memf}},
  bus/.style={blk={busd}{busf}},  per/.style={blk={prd}{prf}},
  sec/.style={blk={secd}{secf}},  pwr/.style={blk={pwrd}{pwrf}},
  med/.style={blk={medd}{medf}},
  comb/.style ={draw=copper, fill=coppert, rounded corners=2pt, thick, align=center,
                font=\scriptsize, inner sep=2pt},                       % combinational logic
  hard/.style ={draw=harde, fill=hardf, thick, dashed, align=center,
                font=\scriptsize, inner sep=2pt},                       % hard macro / off-chip
  sram/.style ={draw=memd, fill=memf, thick, align=center, font=\scriptsize,
                inner sep=2pt},                                         % SRAM array, not flip-flops
  ff/.style   ={draw=#1, fill=#1!30, line width=0.35pt},                % one flip-flop square
  ffgone/.style={draw=muted, fill=white, densely dashed, line width=0.35pt},
  panel/.style={draw=muted, rounded corners=6pt, line width=0.9pt},
  ptitle/.style={anchor=north west, font=\normalsize\bfseries, text=ink, inner sep=0pt},
  note/.style ={font=\scriptsize, text=ink, align=left, inner sep=0pt},
  lab/.style  ={font=\tiny, text=muted, inner sep=1pt},
  axi/.style  ={{Stealth[length=2.6mm]}-{Stealth[length=2.6mm]}, line width=2.2pt, copper},
  axi1/.style ={-{Stealth[length=2.6mm]}, line width=2.2pt, copper},
  apb/.style  ={{Stealth[length=1.8mm]}-{Stealth[length=1.8mm]}, line width=1pt, copper},
  apb1/.style ={-{Stealth[length=1.8mm]}, line width=1pt, copper},
  dat/.style  ={-{Stealth[length=1.8mm]}, thick, pcbmid},
  ctl/.style  ={-{Stealth[length=1.8mm]}, densely dashed, thick, muted},
  irq/.style  ={-{Stealth[length=2mm]}, densely dotted, line width=1.3pt, bad},
  num/.style  ={circle, fill=copper, text=white, font=\scriptsize\bfseries, inner sep=1pt,
                minimum size=0.46cm},
  ours/.style ={draw=bad, line width=1.2pt, densely dashed, rounded corners=4pt},
}

% ---------------- helpers ----------------
\newcommand{\rn}[1]{\texttt{#1}}
% \ffrow{x}{y}{n}{cell}{style}: n flip-flops in a row, bottom-left corner at (x,y)
\newcommand{\ffrow}[5]{\foreach \q in {1,...,#3}{\draw[#5] ({#1+(\q-1)*#4},{#2}) rectangle ++(#4,#4);}}
% \ffcol{x}{y}{n}{cell}{style}: n flip-flops in a column, top-left corner at (x,y)
\newcommand{\ffcol}[5]{\foreach \q in {1,...,#3}{\draw[#5] ({#1},{#2-\q*#4}) rectangle ++(#4,#4);}}
% \ffgrid{x}{ytop}{n}{per row}{cell}{style}: n flip-flops, filled row by row from the top-left
\newcommand{\ffgrid}[6]{\pgfmathtruncatemacro{\gl}{#3-1}%
  \foreach \b in {0,...,\gl}{\pgfmathtruncatemacro{\gc}{mod(\b,#4)}\pgfmathtruncatemacro{\gr}{div(\b,#4)}%
    \draw[#6] ({#1+\gc*#5},{#2-(\gr+1)*#5}) rectangle ++(#5,#5);}}

% every figure: same page size, title, footer
\newcommand{\canvas}{\useasboundingbox (-0.3,-0.2) rectangle (38.6,30.4);}
\newcommand{\figtitle}[3]{%
  \node[num, minimum size=0.9cm, font=\Large\bfseries] at (0.45,29.55) {#1};
  \node[anchor=west, font=\LARGE\bfseries, text=pcb] at (1.15,29.55) {#2};
  \node[anchor=west, font=\small, text=ink] at (0,28.7) {#3};}
\newcommand{\figfoot}[1]{\node[anchor=west, font=\scriptsize, text=muted] at (0,0.25) {#1};}
% \panel{x0}{y0}{x1}{y1}{title}
\newcommand{\panel}[5]{\draw[panel] (#1,#2) rectangle (#3,#4);
  \node[ptitle] at ({#1+0.3},{#4-0.25}) {#5};}

% ---------------- timing diagrams: 1 cycle = \cw cm, low at y, high at y+0.45 ----------------
\newcommand{\cw}{1.0}
\newcommand{\wclk}[3]{% x0 y ncycles
  \foreach \k in {1,...,#3}{\draw[pcb, thick] ({#1+(\k-1)*\cw},#2) -- ++(0,0.45) -- ++(0.5*\cw,0)
     -- ++(0,-0.45) -- ++(0.5*\cw,0);}}
\newcommand{\wsig}[3]{% x0 y {level per cycle, e.g. 0,1,1,0}
  \gdef\wprev{x}%
  \foreach \v [count=\k from 0] in {#3}{%
    \draw[pcb, thick] ({#1+\k*\cw},{#2+0.45*\v}) -- ++(\cw,0);
    \if x\wprev\else\ifnum\wprev=\v\else\draw[pcb, thick] ({#1+\k*\cw},#2) -- ++(0,0.45);\fi\fi
    \global\let\wprev\v}}
\newcommand{\wbus}[5]{% x0 y c0 c1 text: one value held from cycle c0 to cycle c1
  \draw[pcbmid, thick, fill=white] ({#1+(#3)*\cw},{#2+0.225}) -- ({#1+(#3)*\cw+0.08},{#2+0.45})
    -- ({#1+(#4)*\cw-0.08},{#2+0.45}) -- ({#1+(#4)*\cw},{#2+0.225}) -- ({#1+(#4)*\cw-0.08},#2)
    -- ({#1+(#3)*\cw+0.08},#2) -- cycle;
  \node[font=\tiny\ttfamily] at ({#1+((#3)+(#4))*\cw/2},{#2+0.225}) {#5};}
\newcommand{\widle}[4]{% x0 y c0 c1: bus idle (no value)
  \draw[muted, thick] ({#1+(#3)*\cw},{#2+0.225}) -- ({#1+(#4)*\cw},{#2+0.225});}
\newcommand{\wname}[3]{\node[anchor=east, font=\scriptsize\ttfamily] at ({#1-0.15},{#2+0.225}) {#3};}
.
\usepackage[T1]{fontenc}
\usepackage[scaled=0.95]{helvet}
\renewcommand{\familydefault}{\sfdefault}
\usepackage{booktabs}
\usepackage{amsmath}
\usepackage{amssymb}
\usetikzlibrary{arrows.meta,calc,positioning,fit,shapes.geometric,shapes.misc}

% ---------------- palette ----------------
\definecolor{pcb}{HTML}{0B2E33}
\definecolor{pcbmid}{HTML}{1F5C61}
\definecolor{tint}{HTML}{E6F0EF}
\definecolor{copper}{HTML}{B5651D}
\definecolor{coppert}{HTML}{F6E8DC}
\definecolor{ink}{HTML}{1F2A2E}
\definecolor{muted}{HTML}{5B6B70}
\definecolor{hardf}{HTML}{DCE6F2}
\definecolor{harde}{HTML}{2F4A7A}
\definecolor{good}{HTML}{2E7D32}
\definecolor{bad}{HTML}{B3261E}
% one colour per kind of block: border (d) / fill (f)
\definecolor{cpud}{HTML}{2B5C9E} \definecolor{cpuf}{HTML}{D6E4F5}
\definecolor{gpud}{HTML}{2E7D32} \definecolor{gpuf}{HTML}{D5EDD0}
\definecolor{npud}{HTML}{0B2E33} \definecolor{npuf}{HTML}{BFE0DB}
\definecolor{memd}{HTML}{5E3A8C} \definecolor{memf}{HTML}{E4DAF3}
\definecolor{busd}{HTML}{B5651D} \definecolor{busf}{HTML}{F6E8DC}
\definecolor{prd}{HTML}{546E7A}  \definecolor{prf}{HTML}{E3EAED}
\definecolor{secd}{HTML}{B3261E} \definecolor{secf}{HTML}{F7D9D6}
\definecolor{pwrd}{HTML}{8A6800} \definecolor{pwrf}{HTML}{FFF0B8}
\definecolor{medd}{HTML}{8E2463} \definecolor{medf}{HTML}{F5D6E6}
\definecolor{skd}{HTML}{2E7D32}  \definecolor{skf}{HTML}{BFE3B4}   % skew buffers, as in the register figure

\tikzset{
  blk/.style 2 args={draw=#1, fill=#2, rounded corners=4pt, thick, align=center,
                     font=\scriptsize, inner sep=3pt},
  cpu/.style={blk={cpud}{cpuf}},  gpu/.style={blk={gpud}{gpuf}},
  npu/.style={blk={npud}{npuf}},  mem/.style={blk={memd}{memf}},
  bus/.style={blk={busd}{busf}},  per/.style={blk={prd}{prf}},
  sec/.style={blk={secd}{secf}},  pwr/.style={blk={pwrd}{pwrf}},
  med/.style={blk={medd}{medf}},
  comb/.style ={draw=copper, fill=coppert, rounded corners=2pt, thick, align=center,
                font=\scriptsize, inner sep=2pt},                       % combinational logic
  hard/.style ={draw=harde, fill=hardf, thick, dashed, align=center,
                font=\scriptsize, inner sep=2pt},                       % hard macro / off-chip
  sram/.style ={draw=memd, fill=memf, thick, align=center, font=\scriptsize,
                inner sep=2pt},                                         % SRAM array, not flip-flops
  ff/.style   ={draw=#1, fill=#1!30, line width=0.35pt},                % one flip-flop square
  ffgone/.style={draw=muted, fill=white, densely dashed, line width=0.35pt},
  panel/.style={draw=muted, rounded corners=6pt, line width=0.9pt},
  ptitle/.style={anchor=north west, font=\normalsize\bfseries, text=ink, inner sep=0pt},
  note/.style ={font=\scriptsize, text=ink, align=left, inner sep=0pt},
  lab/.style  ={font=\tiny, text=muted, inner sep=1pt},
  axi/.style  ={{Stealth[length=2.6mm]}-{Stealth[length=2.6mm]}, line width=2.2pt, copper},
  axi1/.style ={-{Stealth[length=2.6mm]}, line width=2.2pt, copper},
  apb/.style  ={{Stealth[length=1.8mm]}-{Stealth[length=1.8mm]}, line width=1pt, copper},
  apb1/.style ={-{Stealth[length=1.8mm]}, line width=1pt, copper},
  dat/.style  ={-{Stealth[length=1.8mm]}, thick, pcbmid},
  ctl/.style  ={-{Stealth[length=1.8mm]}, densely dashed, thick, muted},
  irq/.style  ={-{Stealth[length=2mm]}, densely dotted, line width=1.3pt, bad},
  num/.style  ={circle, fill=copper, text=white, font=\scriptsize\bfseries, inner sep=1pt,
                minimum size=0.46cm},
  ours/.style ={draw=bad, line width=1.2pt, densely dashed, rounded corners=4pt},
}

% ---------------- helpers ----------------
\newcommand{\rn}[1]{\texttt{#1}}
% \ffrow{x}{y}{n}{cell}{style}: n flip-flops in a row, bottom-left corner at (x,y)
\newcommand{\ffrow}[5]{\foreach \q in {1,...,#3}{\draw[#5] ({#1+(\q-1)*#4},{#2}) rectangle ++(#4,#4);}}
% \ffcol{x}{y}{n}{cell}{style}: n flip-flops in a column, top-left corner at (x,y)
\newcommand{\ffcol}[5]{\foreach \q in {1,...,#3}{\draw[#5] ({#1},{#2-\q*#4}) rectangle ++(#4,#4);}}
% \ffgrid{x}{ytop}{n}{per row}{cell}{style}: n flip-flops, filled row by row from the top-left
\newcommand{\ffgrid}[6]{\pgfmathtruncatemacro{\gl}{#3-1}%
  \foreach \b in {0,...,\gl}{\pgfmathtruncatemacro{\gc}{mod(\b,#4)}\pgfmathtruncatemacro{\gr}{div(\b,#4)}%
    \draw[#6] ({#1+\gc*#5},{#2-(\gr+1)*#5}) rectangle ++(#5,#5);}}

% every figure: same page size, title, footer
\newcommand{\canvas}{\useasboundingbox (-0.3,-0.2) rectangle (38.6,30.4);}
\newcommand{\figtitle}[3]{%
  \node[num, minimum size=0.9cm, font=\Large\bfseries] at (0.45,29.55) {#1};
  \node[anchor=west, font=\LARGE\bfseries, text=pcb] at (1.15,29.55) {#2};
  \node[anchor=west, font=\small, text=ink] at (0,28.7) {#3};}
\newcommand{\figfoot}[1]{\node[anchor=west, font=\scriptsize, text=muted] at (0,0.25) {#1};}
% \panel{x0}{y0}{x1}{y1}{title}
\newcommand{\panel}[5]{\draw[panel] (#1,#2) rectangle (#3,#4);
  \node[ptitle] at ({#1+0.3},{#4-0.25}) {#5};}

% ---------------- timing diagrams: 1 cycle = \cw cm, low at y, high at y+0.45 ----------------
\newcommand{\cw}{1.0}
\newcommand{\wclk}[3]{% x0 y ncycles
  \foreach \k in {1,...,#3}{\draw[pcb, thick] ({#1+(\k-1)*\cw},#2) -- ++(0,0.45) -- ++(0.5*\cw,0)
     -- ++(0,-0.45) -- ++(0.5*\cw,0);}}
\newcommand{\wsig}[3]{% x0 y {level per cycle, e.g. 0,1,1,0}
  \gdef\wprev{x}%
  \foreach \v [count=\k from 0] in {#3}{%
    \draw[pcb, thick] ({#1+\k*\cw},{#2+0.45*\v}) -- ++(\cw,0);
    \if x\wprev\else\ifnum\wprev=\v\else\draw[pcb, thick] ({#1+\k*\cw},#2) -- ++(0,0.45);\fi\fi
    \global\let\wprev\v}}
\newcommand{\wbus}[5]{% x0 y c0 c1 text: one value held from cycle c0 to cycle c1
  \draw[pcbmid, thick, fill=white] ({#1+(#3)*\cw},{#2+0.225}) -- ({#1+(#3)*\cw+0.08},{#2+0.45})
    -- ({#1+(#4)*\cw-0.08},{#2+0.45}) -- ({#1+(#4)*\cw},{#2+0.225}) -- ({#1+(#4)*\cw-0.08},#2)
    -- ({#1+(#3)*\cw+0.08},#2) -- cycle;
  \node[font=\tiny\ttfamily] at ({#1+((#3)+(#4))*\cw/2},{#2+0.225}) {#5};}
\newcommand{\widle}[4]{% x0 y c0 c1: bus idle (no value)
  \draw[muted, thick] ({#1+(#3)*\cw},{#2+0.225}) -- ({#1+(#4)*\cw},{#2+0.225});}
\newcommand{\wname}[3]{\node[anchor=east, font=\scriptsize\ttfamily] at ({#1-0.15},{#2+0.225}) {#3};}
.
\usepackage[T1]{fontenc}
\usepackage[scaled=0.95]{helvet}
\renewcommand{\familydefault}{\sfdefault}
\usepackage{booktabs}
\usepackage{amsmath}
\usepackage{amssymb}
\usetikzlibrary{arrows.meta,calc,positioning,fit,shapes.geometric,shapes.misc}

% ---------------- palette ----------------
\definecolor{pcb}{HTML}{0B2E33}
\definecolor{pcbmid}{HTML}{1F5C61}
\definecolor{tint}{HTML}{E6F0EF}
\definecolor{copper}{HTML}{B5651D}
\definecolor{coppert}{HTML}{F6E8DC}
\definecolor{ink}{HTML}{1F2A2E}
\definecolor{muted}{HTML}{5B6B70}
\definecolor{hardf}{HTML}{DCE6F2}
\definecolor{harde}{HTML}{2F4A7A}
\definecolor{good}{HTML}{2E7D32}
\definecolor{bad}{HTML}{B3261E}
% one colour per kind of block: border (d) / fill (f)
\definecolor{cpud}{HTML}{2B5C9E} \definecolor{cpuf}{HTML}{D6E4F5}
\definecolor{gpud}{HTML}{2E7D32} \definecolor{gpuf}{HTML}{D5EDD0}
\definecolor{npud}{HTML}{0B2E33} \definecolor{npuf}{HTML}{BFE0DB}
\definecolor{memd}{HTML}{5E3A8C} \definecolor{memf}{HTML}{E4DAF3}
\definecolor{busd}{HTML}{B5651D} \definecolor{busf}{HTML}{F6E8DC}
\definecolor{prd}{HTML}{546E7A}  \definecolor{prf}{HTML}{E3EAED}
\definecolor{secd}{HTML}{B3261E} \definecolor{secf}{HTML}{F7D9D6}
\definecolor{pwrd}{HTML}{8A6800} \definecolor{pwrf}{HTML}{FFF0B8}
\definecolor{medd}{HTML}{8E2463} \definecolor{medf}{HTML}{F5D6E6}
\definecolor{skd}{HTML}{2E7D32}  \definecolor{skf}{HTML}{BFE3B4}   % skew buffers, as in the register figure

\tikzset{
  blk/.style 2 args={draw=#1, fill=#2, rounded corners=4pt, thick, align=center,
                     font=\scriptsize, inner sep=3pt},
  cpu/.style={blk={cpud}{cpuf}},  gpu/.style={blk={gpud}{gpuf}},
  npu/.style={blk={npud}{npuf}},  mem/.style={blk={memd}{memf}},
  bus/.style={blk={busd}{busf}},  per/.style={blk={prd}{prf}},
  sec/.style={blk={secd}{secf}},  pwr/.style={blk={pwrd}{pwrf}},
  med/.style={blk={medd}{medf}},
  comb/.style ={draw=copper, fill=coppert, rounded corners=2pt, thick, align=center,
                font=\scriptsize, inner sep=2pt},                       % combinational logic
  hard/.style ={draw=harde, fill=hardf, thick, dashed, align=center,
                font=\scriptsize, inner sep=2pt},                       % hard macro / off-chip
  sram/.style ={draw=memd, fill=memf, thick, align=center, font=\scriptsize,
                inner sep=2pt},                                         % SRAM array, not flip-flops
  ff/.style   ={draw=#1, fill=#1!30, line width=0.35pt},                % one flip-flop square
  ffgone/.style={draw=muted, fill=white, densely dashed, line width=0.35pt},
  panel/.style={draw=muted, rounded corners=6pt, line width=0.9pt},
  ptitle/.style={anchor=north west, font=\normalsize\bfseries, text=ink, inner sep=0pt},
  note/.style ={font=\scriptsize, text=ink, align=left, inner sep=0pt},
  lab/.style  ={font=\tiny, text=muted, inner sep=1pt},
  axi/.style  ={{Stealth[length=2.6mm]}-{Stealth[length=2.6mm]}, line width=2.2pt, copper},
  axi1/.style ={-{Stealth[length=2.6mm]}, line width=2.2pt, copper},
  apb/.style  ={{Stealth[length=1.8mm]}-{Stealth[length=1.8mm]}, line width=1pt, copper},
  apb1/.style ={-{Stealth[length=1.8mm]}, line width=1pt, copper},
  dat/.style  ={-{Stealth[length=1.8mm]}, thick, pcbmid},
  ctl/.style  ={-{Stealth[length=1.8mm]}, densely dashed, thick, muted},
  irq/.style  ={-{Stealth[length=2mm]}, densely dotted, line width=1.3pt, bad},
  num/.style  ={circle, fill=copper, text=white, font=\scriptsize\bfseries, inner sep=1pt,
                minimum size=0.46cm},
  ours/.style ={draw=bad, line width=1.2pt, densely dashed, rounded corners=4pt},
}

% ---------------- helpers ----------------
\newcommand{\rn}[1]{\texttt{#1}}
% \ffrow{x}{y}{n}{cell}{style}: n flip-flops in a row, bottom-left corner at (x,y)
\newcommand{\ffrow}[5]{\foreach \q in {1,...,#3}{\draw[#5] ({#1+(\q-1)*#4},{#2}) rectangle ++(#4,#4);}}
% \ffcol{x}{y}{n}{cell}{style}: n flip-flops in a column, top-left corner at (x,y)
\newcommand{\ffcol}[5]{\foreach \q in {1,...,#3}{\draw[#5] ({#1},{#2-\q*#4}) rectangle ++(#4,#4);}}
% \ffgrid{x}{ytop}{n}{per row}{cell}{style}: n flip-flops, filled row by row from the top-left
\newcommand{\ffgrid}[6]{\pgfmathtruncatemacro{\gl}{#3-1}%
  \foreach \b in {0,...,\gl}{\pgfmathtruncatemacro{\gc}{mod(\b,#4)}\pgfmathtruncatemacro{\gr}{div(\b,#4)}%
    \draw[#6] ({#1+\gc*#5},{#2-(\gr+1)*#5}) rectangle ++(#5,#5);}}

% every figure: same page size, title, footer
\newcommand{\canvas}{\useasboundingbox (-0.3,-0.2) rectangle (38.6,30.4);}
\newcommand{\figtitle}[3]{%
  \node[num, minimum size=0.9cm, font=\Large\bfseries] at (0.45,29.55) {#1};
  \node[anchor=west, font=\LARGE\bfseries, text=pcb] at (1.15,29.55) {#2};
  \node[anchor=west, font=\small, text=ink] at (0,28.7) {#3};}
\newcommand{\figfoot}[1]{\node[anchor=west, font=\scriptsize, text=muted] at (0,0.25) {#1};}
% \panel{x0}{y0}{x1}{y1}{title}
\newcommand{\panel}[5]{\draw[panel] (#1,#2) rectangle (#3,#4);
  \node[ptitle] at ({#1+0.3},{#4-0.25}) {#5};}

% ---------------- timing diagrams: 1 cycle = \cw cm, low at y, high at y+0.45 ----------------
\newcommand{\cw}{1.0}
\newcommand{\wclk}[3]{% x0 y ncycles
  \foreach \k in {1,...,#3}{\draw[pcb, thick] ({#1+(\k-1)*\cw},#2) -- ++(0,0.45) -- ++(0.5*\cw,0)
     -- ++(0,-0.45) -- ++(0.5*\cw,0);}}
\newcommand{\wsig}[3]{% x0 y {level per cycle, e.g. 0,1,1,0}
  \gdef\wprev{x}%
  \foreach \v [count=\k from 0] in {#3}{%
    \draw[pcb, thick] ({#1+\k*\cw},{#2+0.45*\v}) -- ++(\cw,0);
    \if x\wprev\else\ifnum\wprev=\v\else\draw[pcb, thick] ({#1+\k*\cw},#2) -- ++(0,0.45);\fi\fi
    \global\let\wprev\v}}
\newcommand{\wbus}[5]{% x0 y c0 c1 text: one value held from cycle c0 to cycle c1
  \draw[pcbmid, thick, fill=white] ({#1+(#3)*\cw},{#2+0.225}) -- ({#1+(#3)*\cw+0.08},{#2+0.45})
    -- ({#1+(#4)*\cw-0.08},{#2+0.45}) -- ({#1+(#4)*\cw},{#2+0.225}) -- ({#1+(#4)*\cw-0.08},#2)
    -- ({#1+(#3)*\cw+0.08},#2) -- cycle;
  \node[font=\tiny\ttfamily] at ({#1+((#3)+(#4))*\cw/2},{#2+0.225}) {#5};}
\newcommand{\widle}[4]{% x0 y c0 c1: bus idle (no value)
  \draw[muted, thick] ({#1+(#3)*\cw},{#2+0.225}) -- ({#1+(#4)*\cw},{#2+0.225});}
\newcommand{\wname}[3]{\node[anchor=east, font=\scriptsize\ttfamily] at ({#1-0.15},{#2+0.225}) {#3};}

\begin{document}
\begin{tikzpicture}[line join=round]
\canvas
\figtitle{4}{GPU: thousands of small lanes running the same instruction}
  {A GPU is many shader cores. Each core runs groups of 32 threads (a warp) in lock-step: one instruction, 32 data values.
   Vendor names differ (NVIDIA SM / warp, AMD CU / wavefront); the ideas are the same.}

% =====================================================================
% (A) GPU top level
% =====================================================================
\panel{0}{13.6}{12.5}{28.1}{GPU top level}
\node[gpu, fill=white, minimum width=11.8cm, minimum height=0.9cm, anchor=north west] (hif) at (0.35,27.3)
  {\textbf{host interface}: AXI-Lite registers (doorbell, status) + IRQ wire to the GIC};
\node[gpu, fill=white, minimum width=11.8cm, minimum height=1.2cm, anchor=north west, text width=11.4cm] (cp) at (0.35,26.15)
  {\textbf{command processor}: reads the command buffer the driver wrote in DRAM
   (``draw 2 M triangles'', ``run kernel on 1 M threads'')};
\node[gpu, fill=white, minimum width=11.8cm, minimum height=0.9cm, anchor=north west] (wd) at (0.35,24.7)
  {\textbf{work distributor}: cuts the job into thread blocks and sends them to free cores};
\foreach \r in {0,...,3}{\foreach \c in {0,...,3}{
  \node[gpu, minimum width=2.75cm, minimum height=0.95cm, anchor=north west, font=\tiny]
    at ({0.35+\c*2.95},{23.55-\r*1.1}) {\textbf{shader core \the\numexpr\r*4+\c\relax}\\64 lanes};}}
\node[comb, minimum width=5.75cm, minimum height=1.0cm, anchor=north west, text width=5.5cm, font=\tiny] at (0.35,19.0)
  {\textbf{fixed function} (graphics only): tiler, rasterizer, depth test, blend};
\node[sram, minimum width=5.85cm, minimum height=1.0cm, anchor=north west, font=\tiny] at (6.3,19.0)
  {\textbf{tile buffer}: one screen tile on chip};
\node[sram, minimum width=11.8cm, minimum height=1.0cm, anchor=north west] at (0.35,17.8)
  {\textbf{L2 cache 2 MB}: shared by all cores (SRAM)};
\node[gpu, fill=white, minimum width=11.8cm, minimum height=0.9cm, anchor=north west] at (0.35,16.6)
  {\textbf{GPU MMU} + AXI master ports $\rightarrow$ NoC $\rightarrow$ DRAM};
\node[note, font=\small, anchor=north west, text width=11.8cm] at (0.35,15.4) {%
  16 cores $\times$ 64 lanes = 1,024 FP32 lanes. At 1 GHz with one multiply-add (2 operations) per lane per clock:
  $\approx$2 TFLOPS. Example numbers.};

% =====================================================================
% (B) one shader core
% =====================================================================
\panel{12.8}{13.6}{38.3}{28.1}{Inside one shader core: 4 partitions $\times$ 16 lanes}
\node[gpu, fill=white, minimum width=24.9cm, minimum height=0.7cm, anchor=north west] at (13.1,27.35)
  {\textbf{instruction cache} (SRAM) $\cdot$ shared by the 4 partitions};
\foreach \p in {0,...,3}{
  \pgfmathsetmacro{\x}{13.1+\p*6.3}
  \draw[gpud, thick, rounded corners=4pt] (\x,17.95) rectangle ++(6.0,8.25);
  \node[font=\scriptsize\bfseries, text=gpud, anchor=north west] at (\x+0.1,26.15) {partition \p};
  \node[gpu, fill=white, minimum width=5.7cm, minimum height=0.95cm, anchor=north west, font=\tiny, text width=5.5cm]
    at (\x+0.15,25.7) {\textbf{warp scheduler}: holds up to 16 warps (PC + 32-bit active mask each);
                        picks one ready warp every clock};
  \node[comb, minimum width=5.7cm, minimum height=0.5cm, anchor=north west, font=\tiny] at (\x+0.15,24.6)
    {dispatch: decode, read operands};
  \node[sram, minimum width=5.7cm, minimum height=1.25cm, anchor=north west, font=\tiny, text width=5.5cm]
    at (\x+0.15,23.95) {\textbf{register file 64 KB} = 16,384 $\times$ 32 bit\\
                         SRAM banks; each thread owns up to 255 registers};
  % 16 FP32 lanes
  \foreach \l in {0,...,15}{
    \pgfmathtruncatemacro{\lc}{mod(\l,8)}\pgfmathtruncatemacro{\lr}{div(\l,8)}
    \draw[gpud, fill=gpuf] ({\x+0.15+\lc*0.715},{22.05-\lr*0.5}) rectangle ++(0.62,0.42);
    \node[font=\tiny] at ({\x+0.46+\lc*0.715},{22.26-\lr*0.5}) {FMA};}
  \node[lab, anchor=north west] at (\x+0.1,21.5) {16 FP32 lanes: $a \times b + c$ per clock};
  \node[comb, minimum width=1.75cm, minimum height=0.85cm, anchor=north west, font=\tiny] at (\x+0.15,20.95)
    {4 SFU\\sin, exp, 1/x};
  \node[comb, minimum width=1.75cm, minimum height=0.85cm, anchor=north west, font=\tiny] at (\x+2.1,20.95)
    {4 LD/ST\\addresses};
  % matrix unit drawn as a small MAC grid
  \foreach \mr in {0,...,3}{\foreach \mc in {0,...,3}{
    \draw[ff=npud] ({\x+4.15+\mc*0.42},{20.7-\mr*0.42}) rectangle ++(0.34,0.34);}}
  \node[lab, text=npud, font=\tiny\bfseries, anchor=north] at (\x+4.97,19.25) {matrix unit};
  \node[lab, anchor=north west, text width=3.8cm] at (\x+0.1,19.95) {pipeline registers between every step are flip-flops};
}
\node[sram, minimum width=12.3cm, minimum height=1.25cm, anchor=north west, text width=12.0cm, font=\scriptsize] at (13.1,17.75)
  {\textbf{L1 cache / shared memory 128 KB} (SRAM): threads of one block share data here, like your RAM banks};
\node[gpu, fill=white, minimum width=12.3cm, minimum height=1.25cm, anchor=north west, text width=12.0cm, font=\scriptsize] at (25.7,17.75)
  {\textbf{texture units}: fetch + filter image pixels $\cdot$ path to the GPU L2 cache};
\node[note, font=\small, anchor=north west, text width=24.9cm] at (13.1,16.2) {%
  \textbf{Why SRAM for the register file?} 64 KB = 524,288 bits per partition: far too many for flip-flops.
  The \textcolor{npud}{\textbf{matrix unit}} (``tensor core'') is a small MAC array, the same idea as your 4$\times$4 systolic array:
  it multiplies small int8 / fp16 matrix tiles every clock. Lane counts and sizes are example values.};

% =====================================================================
% (C) SIMT
% =====================================================================
\panel{0}{0.7}{19.0}{13.3}{SIMT: one instruction, 32 threads}
\node[draw=gpud, fill=white, rounded corners=3pt, anchor=north west, font=\scriptsize\ttfamily, align=left] at (0.35,12.55)
  {// kernel: every thread runs this, i = thread number\\y[i] = a * x[i] + y[i];\ \ \ \ \ // one FMA instruction};
\node[anchor=north west, font=\scriptsize\bfseries] at (0.35,11.2) {warp 0 = threads 0--31, same PC:};
\foreach \t in {0,...,31}{
  \draw[gpud, fill=gpuf] ({0.35+\t*0.57},10.0) rectangle ++(0.5,0.6);
  \node[font=\tiny] at ({0.6+\t*0.57},10.3) {\t};}
\node[note, anchor=north west, text width=18.3cm] at (0.35,9.75) {%
  The scheduler issues the FMA \textbf{once}. 16 lanes do threads 0--15 in one clock and threads 16--31 in the next.
  Each thread uses its own registers (its own x[i], y[i]).};
\node[draw=gpud, fill=white, rounded corners=3pt, anchor=north west, font=\scriptsize\ttfamily, align=left] at (0.35,8.3)
  {if (x[i] > 0) y[i] = x[i];\ \ \ \ \ // path A\\else\ \ \ \ \ \ \ \ \ \ y[i] = 0;\ \ \ \ \ \ \ \ // path B};
\node[anchor=west, font=\scriptsize\bfseries] at (0.35,6.6) {mask, path A};
\node[anchor=west, font=\scriptsize\bfseries] at (0.35,5.9) {mask, path B};
\foreach \t in {0,...,31}{
  \pgfmathtruncatemacro{\on}{mod(\t,3)==0 ? 0 : 1}
  \ifnum\on=1
    \draw[gpud, fill=gpud!60] ({3.0+\t*0.49},6.42) rectangle ++(0.42,0.36);
    \draw[muted, fill=white] ({3.0+\t*0.49},5.72) rectangle ++(0.42,0.36);
  \else
    \draw[muted, fill=white] ({3.0+\t*0.49},6.42) rectangle ++(0.42,0.36);
    \draw[bad, fill=bad!45] ({3.0+\t*0.49},5.72) rectangle ++(0.42,0.36);
  \fi}
\node[note, anchor=north west, text width=18.3cm] at (0.35,5.35) {%
  \textbf{Divergence}: threads of one warp take different paths. The warp runs path A with only the A-threads switched on
  (active mask), then path B with the others. Both paths cost time, so branches inside a warp are slow.\\[3pt]
  \textbf{Lesson for AI}: matrix multiply has no branches: every thread does the same multiply-adds.
  That is why GPUs (and NPUs) are good at it.};

% =====================================================================
% (D) latency hiding
% =====================================================================
\panel{19.3}{6.9}{38.3}{13.3}{Hiding memory latency: switch warps every clock}
\foreach \w in {0,...,3}{\node[font=\scriptsize\bfseries, anchor=east] at (21.0,{11.85-\w*0.85}) {warp \w};}
\foreach \w/\s in {0/0, 1/1, 2/2, 3/3}{
  \pgfmathsetmacro{\yy}{11.6-\w*0.85}
  \fill[gpud!60] ({21.2+\s*0.5},\yy) rectangle ++(0.5,0.5);
  \fill[gray!25] ({21.7+\s*0.5},\yy) rectangle ++(6.0,0.5);
  \node[font=\tiny, text=muted] at ({24.7+\s*0.5},{\yy+0.25}) {waiting $\approx$400 cycles for DRAM};
  \fill[gpud!60] ({27.7+\s*0.5},\yy) rectangle ++(0.5,0.5);}
\draw[-{Stealth}, thick] (21.2,8.15) -- (30.0,8.15);
\node[lab, anchor=north east] at (30.0,8.1) {time};
\node[note, font=\small, anchor=north west, text width=7.3cm] at (30.6,12.0) {%
  \textcolor{gpud}{\textbf{green}} = issuing instructions\\
  grey = waiting for memory\\[3pt]
  With 16 warps per scheduler there is almost always one ready, so the lanes stay busy.};
\node[note, anchor=north west, text width=18.3cm] at (19.6,7.85) {%
  CPU hides memory latency with big caches and out-of-order execution. GPU hides it with \textbf{many warps}.};

% =====================================================================
% (E) comparison
% =====================================================================
\panel{19.3}{0.7}{38.3}{6.6}{CPU vs GPU vs NPU}
\node[anchor=north west, font=\footnotesize, inner sep=0pt] at (19.6,5.85) {%
  \renewcommand{\arraystretch}{1.1}%
  \begin{tabular}{@{}l l l l@{}}
    \toprule
                    & \textbf{CPU}                 & \textbf{GPU}                     & \textbf{NPU}\\
    \midrule
    compute units   & 4--12 big cores               & 1,000s of simple lanes           & 1 big MAC array\\
    control         & each core its own PC          & 1 PC per warp of 32              & fixed dataflow, 1 FSM\\
    best at         & branches, OS, any code        & graphics, fp32/fp16 math, training & int8 inference\\
    energy per op   & highest                       & middle                           & lowest\\
    on your board   & (next: RISC-V)                & --                               & \rn{mm\_core} (4$\times$4)\\
    \bottomrule
  \end{tabular}};

\figfoot{Lane counts, register-file sizes and clock rates are example values for a recent GPU, not one specific product.}
\end{tikzpicture}
\end{document}
